\section{Version B: an exogenous material intensity of capital}\label{sec:sp:B}

This section develops the planner's problem in full under Version B, the closure in which the material
intensity of new capital goods, $\Phi^I(t)$, is exogenous. It is written to be read on its own: the
problem is stated, the optimality conditions are derived, and their content is worked out without
reference to Section~\ref{sec:sp:A}, which does the same for Version A. Only the setup of
Section~\ref{sec:sp:setup} is presupposed. Section~\ref{sec:sp:B:compare} then collects the differences
between the two versions in one place, and is the only part of this section that requires
Section~\ref{sec:sp:A} to have been read.


\subsection{The closure and what it determines}\label{sec:sp:B:closure}

Under Version B the material intensity of new capital goods is a given function of time, $\Phi^I(t)$,
independent of the material intensity of final-good activity in general. It has a direct technological
reading: $\Phi^I$ is the mass required to build a unit of capital goods out of the final good. Three
consequences shape the problem below.

\smalltitle{(i) The accounting is recursive.} The two aggregate balances
\eqref{eq:sp:setup:ybalance_generic} and~\eqref{eq:sp:setup:ybalance_Omega} determine the level $\Omega$
and the material intensity of the final good $\Phi^Y$ pointwise, from~\eqref{eq:sp:setup:Omega},
\begin{align}
\Omega &= \frac{R-\Phi^IG}{\mathcal D},
&
\Phi^Y &= \Omega\,\frac{\mathcal D-\omega^YY}{Y}+\Phi^I\frac{G}{Y},
&
\mathcal D &= \omega^YY+\omega^{N}C^N+\omega^{D}C^D+\omega^WC^W+\omega^CC .
\label{eq:sp:B:Omega}
\end{align}
Neither relation restricts the allocation: given any feasible path, the two compute $\Omega$ and
$\Phi^Y$ after the fact. This is what makes the accounting recursive --- the allocation can be solved
first and the intensities read off afterwards --- and it is the single structural fact from which most
of what follows descends.

\smalltitle{(ii) One restriction survives.} What~\eqref{eq:sp:B:Omega} does impose is $\Omega\ge0$,
that is $\Phi^IG\le R$: the economy cannot embody more mass in new capital than entered it. With
$\Phi^I$ exogenous this is a genuine restriction on the feasible set, and it is carried into the
problem as an \emph{inequality} with a complementary-slackness multiplier. It is not an equality
whose multiplier is nonzero by construction, which is what imposing the materials balance as a
constraint would amount to. Note that $\Phi^Y\ge0$ is an implication rather than a separate condition,
since every term in the second expression in~\eqref{eq:sp:B:Omega} is non-negative once $\Omega$ is.

\smalltitle{(iii) The ledger is linear in $W$.} Because $\Phi^I$ does not depend on $W$, neither does
$\dot M^K=\Phi^IG-\delta M^K$, and the collected ledger~\eqref{eq:sp:setup:ledger-collected} is already
a solution for the waste flow, $W=\left(N-\dot M^K+\mathcal N\right)/\left(1-a\varpi\right)$. The
relative intensities $\omega^j$ appear neither in it nor in the restriction $\Phi^IG\le R$, since the
positive base $\mathcal D$ divides out of the latter.


\subsection{The planner's problem}\label{sec:sp:B:problem}

\smalltitle{Controls and states.} The planner chooses paths for the seven controls
\begin{align}
\left\{C,\;I,\;N,\;D,\;\varpi,\;K^R,\;W\right\},
\label{eq:sp:B:controls}
\end{align}
which determine the remaining flows through the definitions
\begin{align}
R^R &= \mathcal R\left(\varpi W,K^R\right),
&
R &= N+R^R,
&
Y &= F\left(K-K^R,\,R,\,P\right),
&
G &= I+\delta K,
\label{eq:sp:B:definitions}
\end{align}
and it carries the five states
\begin{align}
\left\{K,\;S,\;X,\;P,\;M^K\right\},
\label{eq:sp:B:states}
\end{align}
with $K_0,S_0,X_0,P_0,M^K_0$ given. Note that $\Omega$ is \emph{not} among the controls: by~(i) above
it is determined residually and does not feed back, so it can be left out of the problem entirely and
recovered afterwards from~\eqref{eq:sp:B:Omega}.

\smalltitle{Why $W$ is listed as a control.} The ledger determines $W$ residually too, so it is not an
independent object either; listing it among the controls and imposing the ledger as an equality
constraint is the equivalent formulation that keeps the shadow price of waste visible. Substituting the
solved ledger and dropping $W$ from~\eqref{eq:sp:B:controls} yields the same allocation, but the
multiplier $p^W$ --- the object every environmental instrument in the market economy is built from ---
would then have to be recovered afterwards. Remark~\ref{rem:sp:B:reduced} states the equivalence.

\begin{problem}[Planner, Version B]\label{prob:sp:B}
Choose $\left\{C,I,N,D,\varpi,K^R,W\right\}_{t\ge0}$ to maximise
$U_0=\int_0^\infty\left[u(C)-v(P)\right]e^{-\rho t}dt$ subject to, at every $t$,
\begin{align}
&\text{\emph{goods:}}
&
F\left(K-K^R,R,P\right) &= C+I+\delta K+C^N(N,S)+C^D(D,X)+c^W\!\left(\varpi\right)W,
\label{eq:sp:B:goods}
\\[0.2em]
&\text{\emph{ledger:}}
&
W &= R-\Phi^I G+\delta M^K+\mathcal N,
\label{eq:sp:B:ledger}
\\[0.2em]
&\text{\emph{mass:}}
&
\Phi^I G &\le R,
\label{eq:sp:B:mass}
\end{align}
and to the laws of motion
\begin{align}
\dot K &= I,
&
\dot S &= D-N,
&
\dot X &= D,
&
\dot P &= \Xi-\theta(P)P,
&
\dot M^K &= \Phi^I G-\delta M^K,
\label{eq:sp:B:states-motion}
\end{align}
with $\mathcal N=\Omega^{N,S}N+\Omega^{D,S}D$, $c^W(\varpi)=c^c+c^T(\varpi)$, the definitions
in~\eqref{eq:sp:B:definitions}, the emission flow
\begin{align}
\Xi &= \left(1-\varpi+d^W\varpi\right)W-d^WR^R,
\label{eq:sp:B:Xi}
\end{align}
and $C,I+\delta K,N,D,K^R,W\ge0$, $\varpi\in\left[0,1\right]$.
\end{problem}

Three features of this statement deserve marking before any algebra.

First, \eqref{eq:sp:B:ledger} is an \emph{identity}, not a technology. It says where the mass went, and
it holds along every conceivable path; what makes it bite is that it ties the waste flow $W$ --- which
is costly to collect, costly to treat and damaging when released --- to the material flow $R$ that
production requires. Nothing in it can be relaxed by choosing differently; it can only be made less
costly by choosing a path on which less mass has to be moved.

Second, \eqref{eq:sp:B:mass} is the only part of the materials accounting that restricts the feasible
set, and it does so as an inequality. This is the point on which the present formulation departs most
sharply from the reduced-form treatment: imposing the materials balance as an \emph{equality}
constraint with a freely signed multiplier would attach a nonzero shadow price to mass conservation by
construction, and that price would show up as a wedge on output and on consumption. Here the multiplier
on~\eqref{eq:sp:B:mass} is zero except at the corner where every tonne entering production is embodied
in new capital.

Third, $\Xi$ is written in the form~\eqref{eq:sp:B:Xi}, linear in $W$ and $R^R$ separately, so that no
derivative of $a(\cdot)$ ever appears inside a pollution term. All the curvature of the recycling
technology enters through $\mathcal R$ and through $\alpha$ and $a'$ alone; recall
from~\eqref{eq:sp:setup:recycling-derivatives} that $\mathcal R_T=\alpha\equiv a-a'x$ is the marginal
recycling yield and $\mathcal R_{K^R}=a'$.

\smalltitle{The Lagrangian.} Write the current-value Hamiltonian with the static constraints adjoined.
It is convenient to normalise every multiplier by $\lambda$, the multiplier on the goods
constraint~\eqref{eq:sp:B:goods}, so that all shadow prices below are measured in units of the final
good rather than in utils. Accordingly, let the current-value costates be
\begin{align}
\underbrace{\lambda q}_{K},
&&
\underbrace{\lambda p^S}_{S},
&&
\underbrace{\lambda p^X}_{X},
&&
\underbrace{-\lambda p^P}_{P},
&&
\underbrace{-\lambda p^M}_{M^K},
\label{eq:sp:B:costates-def}
\end{align}
the two minus signs anticipating that both $P$ and $M^K$ are liabilities, so that $p^P$ and $p^M$ read
as shadow \emph{costs}; the signs are verified in Section~\ref{sec:sp:B:signs}, not assumed. Let
$\lambda p^W$ be the multiplier on the ledger~\eqref{eq:sp:B:ledger} written with $W$ on the left, and
$\lambda\eta\ge0$ the multiplier on~\eqref{eq:sp:B:mass}. Then
\begin{align}
\mathcal L
={}& u(C)-v(P)
\notag\\
&+\lambda\Big[F\left(K-K^R,R,P\right)-C-I-\delta K-C^N(N,S)-C^D(D,X)-c^W\!\left(\varpi\right)W\Big]
\notag\\
&+\lambda q\,I
\;+\;\lambda p^S\left(D-N\right)
\;+\;\lambda p^X D
\;-\;\lambda p^P\left[\Xi-\theta(P)P\right]
\;-\;\lambda p^M\left[\Phi^IG-\delta M^K\right]
\notag\\
&+\lambda p^W\Big[W-N-R^R+\Phi^IG-\delta M^K-\mathcal N\Big]
\;+\;\lambda\eta\Big[N+R^R-\Phi^IG\Big],
\label{eq:sp:B:lagrangian}
\end{align}
with $R=N+R^R$, $R^R=\mathcal R\left(\varpi W,K^R\right)$ and $G=I+\delta K$ substituted throughout.
Note what is absent: no term in $\mathcal D$, and hence no relative intensity $\omega^j$, appears
anywhere in~\eqref{eq:sp:B:lagrangian}.

The sign convention on $p^W$ is worth stating explicitly, because it is the price everything else is
written in terms of. Perturb the ledger to $W=R-\Phi^IG+\delta M^K+\mathcal N-\kappa$, that is, suppose
$\kappa$ tonnes of waste could be made to disappear at no cost. Then
$\partial\mathcal L^*/\partial\kappa=\lambda p^W$, so $p^W>0$ says exactly that waste is socially
costly, and $p^W$ measures how costly, per tonne, in final-good units.


\subsection{First-order conditions}\label{sec:sp:B:foc}

Differentiating~\eqref{eq:sp:B:lagrangian} with respect to each of the seven controls, and dividing
throughout by $\lambda$, gives the following system. Two composite objects recur often enough to be
named at the outset:
\begin{align}
B^{\mathrm B} &\equiv F_R-p^W+d^Wp^P+\eta,
&
p^{N,\mathrm B} &\equiv C^N_N+p^S+\Omega^{N,S}p^W .
\label{eq:sp:B:Bvalue}
\end{align}
$B^{\mathrm B}$ is the net social value of one tonne of \emph{secondary} material: what it produces, less the waste
it will itself become, plus the emission its recovery avoids, plus the mass constraint it relaxes.
$p^{N,\mathrm B}$ is the full marginal social cost of one tonne of \emph{virgin} material delivered to production:
the resources spent extracting it, the rent on the reserve it depletes, and the waste cost of the
overburden its extraction displaces. Section~\ref{sec:sp:B:materials} argues that the relation between
these two objects is the single most informative equation in the model.

\smalltitle{Consumption.}
\begin{align}
u'\!\left(C\right) &= \lambda .
\label{eq:sp:B:foc:C}
\end{align}
No wedge: at the optimum the marginal utility of consumption equals the shadow value of a unit of the
final good, exactly as in an economy with no materials at all.

\smalltitle{Investment.}
\begin{align}
q &= 1+\Phi^I\left(p^M-p^W+\eta\right).
\label{eq:sp:B:foc:I}
\end{align}
The shadow price of installed capital is the goods it costs, \emph{less} the waste that embodying
$\Phi^I$ tonnes in it removes from today's waste stream, \emph{plus} the discounted waste those same
tonnes will become when the capital depreciates. Capital is a temporary materials sink, and $q$ prices
the net value of that service.

\smalltitle{Extraction.}
\begin{align}
F_R+\eta-p^W &= C^N_N+p^S+\Omega^{N,S}p^W,
\qquad\text{equivalently}\qquad
B^{\mathrm B} &= p^{N,\mathrm B}+d^Wp^P .
\label{eq:sp:B:foc:N}
\end{align}

\smalltitle{Exploration.}
\begin{align}
p^S+p^X &= C^D_D+\Omega^{D,S}p^W .
\label{eq:sp:B:foc:D}
\end{align}

\smalltitle{Treatment share.}
\begin{align}
\alpha B^{\mathrm B}+\left(1-d^W\right)p^P &= \frac{dc^T}{d\varpi},
\label{eq:sp:B:foc:varpi}
\end{align}
for interior $\varpi\in\left(0,1\right)$, with the usual inequalities at the corners.

\smalltitle{Recycling capital.}
\begin{align}
F_K &= a'\!\left(x\right)B^{\mathrm B} .
\label{eq:sp:B:foc:KR}
\end{align}

\smalltitle{Waste.}
\begin{align}
p^W &= c^W\!\left(\varpi\right)+\left(1-\varpi+d^W\varpi\right)p^P-\varpi\alpha B^{\mathrm B} .
\label{eq:sp:B:foc:W}
\end{align}

\smalltitle{Complementary slackness.}
\begin{align}
\eta &\ge0,
&
\Phi^IG &\le R,
&
\eta\left[R-\Phi^IG\right] &= 0 .
\label{eq:sp:B:foc:cs}
\end{align}

Two of these are worth expanding on immediately, because their form is not obvious.

\eqref{eq:sp:B:foc:W} defines $p^W$ implicitly, since $B^{\mathrm B}$ contains $p^W$. Solving out,
using~\eqref{eq:sp:B:foc:N} to replace $B^{\mathrm B}$, the shadow cost of a tonne of gross waste satisfies
\begin{align}
p^W\left(1+\Omega^{N,S}\varpi\alpha\right)
&= c^W\!\left(\varpi\right)+\left(1-\varpi+d^W\varpi\right)p^P
-\varpi\alpha\left(C^N_N+p^S+d^Wp^P\right).
\label{eq:sp:B:foc:W-solved}
\end{align}
A tonne of waste costs what it costs to collect and treat, plus the damage from the part of it that
reaches the environment, minus the virgin extraction it displaces at the margin. \emph{Nothing
in~\eqref{eq:sp:B:foc:W-solved} involves any $\omega^j$ or $\Omega$}: the shadow price of waste is a
function of extraction costs, resource rents, damages and the recycling technology alone, and unlike
its Version~A counterpart it is an explicit formula rather than an implicit determination.

\eqref{eq:sp:B:foc:varpi} is the treatment margin. Raising $\varpi$ by one point costs $dc^T/d\varpi$
per tonne of waste and buys two things: $\alpha$ tonnes of secondary material worth $B^{\mathrm B}$ each, and a
reduction of the emitted fraction from $1$ to $d^W$ on the tonne now treated. Since
$c^T(0)=\left.dc^T/d\varpi\right|_{0}=0$ by~\eqref{eq:sp:setup:waste-cost}, the right-hand side vanishes
at $\varpi=0$ while the left-hand side is strictly positive whenever $p^P>0$; so \emph{some} treatment
is always optimal, and the $\varpi=0$ corner is never reached. The upper corner $\varpi=1$ is reached
iff $\alpha B^{\mathrm B}+\left(1-d^W\right)p^P\ge\left.dc^T/d\varpi\right|_{1}$.


\subsection{Costate equations}\label{sec:sp:B:costates}

Each costate obeys $\dot\nu=\rho\nu-\partial\mathcal L/\partial(\text{state})$ with $\nu$ the
current-value costate in~\eqref{eq:sp:B:costates-def}. Because the costates carry the factor $\lambda$,
each such equation produces a term $\dot\lambda/\lambda$, and it is worth defining the resulting
discount rate once:
\begin{align}
r &\equiv \rho-\frac{\dot\lambda}{\lambda}.
\label{eq:sp:B:r}
\end{align}
Since $\lambda=u'(C)$ by~\eqref{eq:sp:B:foc:C}, $r$ is the familiar consumption rate of interest,
$r=\rho+\sigma(C)\dot C/C$ with $\sigma\equiv-Cu''/u'$; and because every shadow price has been
normalised by $\lambda$, \emph{every} price below discounts at this same $r$. That uniformity is not
automatic --- it is a consequence of the normalisation together with $\lambda=u'(C)$, and it is one of
the two places where Version B is materially simpler than Version A.

\smalltitle{Capital.} Using $\partial\mathcal L/\partial K=\lambda\left[F_K-\delta-\delta\Phi^I
\left(p^M-p^W+\eta\right)\right]$ and~\eqref{eq:sp:B:foc:I} to collapse the bracket to $F_K-\delta q$,
\begin{align}
r &= \frac{F_K+\dot q}{q}-\delta .
\label{eq:sp:B:costate:K}
\end{align}
Depreciation is charged at $q$, not at $1$: what wears out is a unit of installed capital, whose social
value is $q$, not a unit of the final good. Together with~\eqref{eq:sp:B:r} this is the consumption
Euler equation,
\begin{align}
\sigma\left(C\right)\frac{\dot C}{C} &= \frac{F_K+\dot q}{q}-\delta-\rho .
\label{eq:sp:B:euler}
\end{align}

\smalltitle{Reserves.} With $\partial\mathcal L/\partial S=-\lambda C^N_S$,
\begin{align}
\dot p^S &= r\,p^S+C^N_S,
\qquad\text{so}\qquad
\frac{\dot p^S}{p^S} = r+\frac{C^N_S}{p^S}<r .
\label{eq:sp:B:costate:S}
\end{align}
This is the modified Hotelling rule: the in-situ rent grows more slowly than the interest rate, by
exactly the marginal cost saving $-C^N_S>0$ that a larger known reserve confers on extraction.

\smalltitle{Discoveries.} With $\partial\mathcal L/\partial X=-\lambda C^D_X$,
\begin{align}
\dot p^X &= r\,p^X+C^D_X,
\qquad\text{so}\qquad
p^X\left(t\right) = -\int_t^\infty C^D_X\!\left(s\right)e^{-\int_t^s r\,du}\,ds<0 .
\label{eq:sp:B:costate:X}
\end{align}
Cumulative discovery is a bad: each unit found today raises the cost of finding the next one forever
after, and $-p^X$ is the present value of that penalty.

\smalltitle{Pollution.} Write the marginal damage of the stock, in final-good units, as
\begin{align}
d^{\mathrm B} &\equiv \frac{v'\!\left(P\right)}{u'\!\left(C\right)}-F_P\;>\;0,
\label{eq:sp:B:damage}
\end{align}
the sum of the direct disutility and the output loss. Then
\begin{align}
\dot p^P &= \Big[r+\theta\!\left(P\right)+\theta'\!\left(P\right)P\Big]p^P-d^{\mathrm B},
\qquad\text{so}\qquad
p^P\left(t\right)=\int_t^\infty d^{\mathrm B}\left(s\right)e^{-\int_t^s\left[r+\theta+\theta'P\right]du}\,ds .
\label{eq:sp:B:costate:P}
\end{align}
The effective discount rate on damages is the interest rate plus the marginal decay rate
$\partial\left[\theta(P)P\right]/\partial P=\theta+\theta'P$. Because $\theta'<0$
by~\eqref{eq:sp:setup:pollution-motion}, sink saturation lowers that rate below the average decay rate
$\theta$, and $p^P$ is correspondingly larger than a naive $d^{\mathrm B}/\left(r+\theta\right)$ calculation would
suggest. Stability of the pollution block requires $\theta+\theta'P>0$; this is an assumption on
$\theta$, and it is maintained throughout.

\smalltitle{Embodied material.} With $\partial\mathcal L/\partial M^K=\lambda\delta\left(p^M-p^W\right)$,
\begin{align}
\dot p^M &= \left(r+\delta\right)p^M-\delta p^W,
\qquad\text{so}\qquad
p^M\left(t\right)=\int_t^\infty \delta p^W\!\left(s\right)
e^{-\left(\int_t^s r\,du\right)-\delta\left(s-t\right)}ds .
\label{eq:sp:B:costate:M}
\end{align}
$p^M$ is the present value of the waste that the currently embodied stock will release as it
depreciates, discounted at $r+\delta$ because a fraction $\delta$ of the liability is discharged each
instant. It is positive whenever $p^W$ is, confirming the sign anticipated
in~\eqref{eq:sp:B:costates-def}, and it is strictly below the contemporaneous $p^W$ whenever $r>0$ ---
postponing a tonne of waste is worth something.

\smalltitle{Transversality.} The problem is closed by
\begin{align}
\lim_{t\to\infty}e^{-\rho t}\lambda q K &= 0,
&
\lim_{t\to\infty}e^{-\rho t}\lambda p^S S &= 0,
&
\lim_{t\to\infty}e^{-\rho t}\lambda p^X X &= 0,
\notag\\
\lim_{t\to\infty}e^{-\rho t}\lambda p^P P &= 0,
&
\lim_{t\to\infty}e^{-\rho t}\lambda p^M M^K &= 0 .
&&
\label{eq:sp:B:tvc}
\end{align}


\subsection{The pricing of materials}\label{sec:sp:B:materials}

Everything in the materials block follows from one equation. Rearranging the extraction
condition~\eqref{eq:sp:B:foc:N},
\begin{align}
\boxed{\;B^{\mathrm B}\;=\;p^{N,\mathrm B}+d^Wp^P\;}
\qquad\text{that is}\qquad
\underbrace{F_R-p^W+d^Wp^P+\eta}_{\text{value of a secondary tonne}}
\;=\;
\underbrace{C^N_N+p^S+\Omega^{N,S}p^W}_{\text{cost of a virgin tonne}}
\;+\;d^Wp^P .
\label{eq:sp:B:arbitrage}
\end{align}
Read from the left, this is a definition of $B^{\mathrm B}$ that has been shown to equal something; read from the
right, it is the materials arbitrage condition. A tonne of material can be obtained in two ways.
Extracting it costs $C^N_N$ in resources, $p^S$ in foregone reserve value, and $\Omega^{N,S}p^W$ in the
waste cost of the overburden it drags up. Recovering it from waste instead avoids all three, and
additionally avoids the emission $d^W$ that the tonne would have caused had it stayed untreated ---
which is why $d^Wp^P$ appears on both sides and cancels out of the comparison of \emph{sources} while
remaining in the valuation of the recovered tonne.

The cancellation is the sharpest check available on the accounting. Note what has dropped out: $p^W$,
the shadow price of waste, appears on both sides of the arbitrage and is not what determines whether
recycling is worthwhile. It is worth being precise about why. A tonne of secondary material does not
\emph{remove} a tonne of waste from the economy; it postpones it. The tonne is recovered, used in
production, and returns as waste on the next pass. What recycling saves is not the waste cost but the
\emph{extraction} cost --- the resources, the rent, and the displaced overburden. This is the
substantive content of treating the materials balance as a ledger: a reduced-form model in which
recycling is credited with avoiding waste double-counts, because the waste is not avoided, only
deferred by one circuit.

Substituting~\eqref{eq:sp:B:arbitrage} back into the definition of $B^{\mathrm B}$ gives the companion statement
for the production margin,
\begin{align}
F_R &= p^{N,\mathrm B}+p^W-\eta .
\label{eq:sp:B:FR}
\end{align}
The marginal product of materials in production covers the full cost of delivering a virgin tonne
\emph{plus} the waste that tonne will eventually become. Here $p^W$ does appear, and it must: a tonne
drawn into the economy leaves as waste exactly once, whatever route it takes on the way, and that is
the point at which the waste cost is charged.


\subsection{What the accounting does not price}\label{sec:sp:B:nowedge}

Two absences in Section~\ref{sec:sp:B:foc} carry as much information as any of the equations.

\begin{proposition}[No wedge on output or consumption]\label{prop:sp:B:nowedge}
Under Version B the consumption condition is $u'(C)=\lambda$, with $\lambda$ the multiplier on the goods
constraint, and the marginal social value of a unit of output is $\lambda$. Neither carries a correction
for the waste that consumption or production generates.
\end{proposition}

This is not the statement that consumption and output generate no waste --- by \eqref{eq:sp:setup:wc}
they generate $\Omega^CC$ and $\Omega^YY$ tonnes of it. It is the statement that at the margin the
planner should not act on that fact. The reason is~\eqref{eq:sp:B:FR}: the mass that becomes
$\Omega^CC$ was drawn into the economy as part of $R$ and was charged $p^W$ at that point. Charging it
again when it is spent on consumption would price the same tonne twice.

The contrast with a reduced-form specification is worth making explicit, because it is the main
methodological payoff of Section~\ref{sec:sp:setup:materialaccounting}. Suppose the
ledger~\eqref{eq:sp:B:ledger} were replaced by a posited waste function with \emph{fixed} intensities,
\begin{align}
W &= \Omega\left(\omega^YY+\omega^CC+\omega^{N}C^N+\omega^{D}C^D+\omega^WC^W\right)+\Phi^K\delta K,
&&\Omega \text{ a parameter},
\label{eq:sp:B:reducedform}
\end{align}
which is~\eqref{eq:sp:setup:totalwaste_generic} with the substitution
\eqref{eq:sp:setup:ybalance_generic2} \emph{not} carried out. Redoing Section~\ref{sec:sp:B:foc} with
\eqref{eq:sp:B:reducedform} in place of the ledger gives, in place of~\eqref{eq:sp:B:foc:C},
\begin{align}
u'\!\left(C\right) &= \lambda\left(1+\Omega^Cp^W\right),
&
\Omega^C&\equiv\Omega\omega^C,
\label{eq:sp:B:spurious-C}
\end{align}
and a parallel wedge $\left(1-\Omega^Yp^W\right)$ on the marginal value of output --- a consumption tax
and an output tax, each proportional to the assumed intensity. These wedges are artefacts. They appear
because \eqref{eq:sp:B:reducedform} lets the planner reduce economy-wide waste by consuming less
\emph{without reducing the mass drawn in}, which conservation of mass does not permit. Once the
accounting is done process by process, $\Omega$ is endogenous, the substitution that produces the ledger
cancels every $\omega^j$, and the wedges vanish. The difference is visible in the optimal policy, not
merely in the derivation: it is the difference between a model that recommends taxing consumption on
environmental grounds and one that does not.

\begin{proposition}[Irrelevance of the relative intensities]\label{prop:sp:B:irrelevance}
Under Version B, no $\omega^j$ and no $\Omega$ appears in any first-order
condition~\eqref{eq:sp:B:foc:C}--\eqref{eq:sp:B:foc:cs}, in any costate
equation~\eqref{eq:sp:B:costate:K}--\eqref{eq:sp:B:costate:M}, or in the constraint set of
Problem~\ref{prob:sp:B}. The optimal allocation is therefore invariant to the relative intensities, and
$\Omega$ and $\Phi^Y$ are recovered afterwards from~\eqref{eq:sp:B:Omega} as reported quantities.
\end{proposition}

The proof is inspection, and the mechanism is the cancellation described
after~\eqref{eq:sp:setup:ledger}: conservation of mass pins total waste down exactly, because every
gram of the accounted flow passed through $R$ on its way in, and how that gram was distributed across
activities on the way through does not change the total. The absolute intensities $\Omega^{N,S}$ and
$\Omega^{D,S}$ are a different matter --- they multiply flows that never passed through $R$, they
appear in~\eqref{eq:sp:B:foc:N} and~\eqref{eq:sp:B:foc:D}, and they do affect the allocation.

Proposition~\ref{prop:sp:B:irrelevance} has an uncomfortable corollary for the empirical side of the
project. Material-flow accounts report intensities, and it is natural to want to discipline the
$\omega^j$ against them; but under Version B no observation of the \emph{allocation} --- no transition
date, no recycling rate, no extraction path --- carries information about them. They are identified only
from the reported mass flows themselves, and they inform nothing else. Version A reverses this, which is
the substantive reason the two closures are both carried rather than one being chosen on grounds of
plausibility; see Section~\ref{sec:sp:B:compare}.


\subsection{Capital as a materials sink, and the signs of the prices}\label{sec:sp:B:signs}

The investment condition~\eqref{eq:sp:B:foc:I} is the one place where the stock-flow structure of
embodied material shows up in a price. Setting $\eta=0$ (the generic case, see below),
\begin{align}
q &= 1-\Phi^I\left(p^W-p^M\right).
\label{eq:sp:B:q}
\end{align}
Building a unit of capital costs one unit of the final good and takes $\Phi^I$ tonnes out of the current
waste stream, worth $\Phi^Ip^W$; those tonnes return as waste over the life of the asset, at a present
value of $\Phi^Ip^M$. Capital is a materials sink, but a temporary one, and $q$ prices the difference.
Since~\eqref{eq:sp:B:costate:M} makes $p^M$ a discounted average of future $p^W$, the bracket is
positive whenever $p^W$ is positive and not falling fast, so $q<1$: the social cost of installed capital
is \emph{below} its goods cost. Along a path on which $p^W$ is constant, \eqref{eq:sp:B:costate:M} gives
$p^M=\delta p^W/\left(r+\delta\right)$ and
\begin{align}
q &= 1-\Phi^Ip^W\frac{r}{r+\delta},
\label{eq:sp:B:q-stationary}
\end{align}
so the wedge is largest for durable capital ($\delta$ small) and vanishes as $\delta\to\infty$, when
capital stores nothing. It also vanishes as $r\to0$, when postponement is worthless.

\smalltitle{The sign of $p^W$.} Nothing guarantees $p^W>0$. From~\eqref{eq:sp:B:foc:W-solved},
\begin{align}
p^W<0
\qquad\Longleftrightarrow\qquad
\varpi\alpha\left(C^N_N+p^S+d^Wp^P\right)
\;>\;
c^W\!\left(\varpi\right)+\left(1-\varpi+d^W\varpi\right)p^P,
\label{eq:sp:B:pW-sign}
\end{align}
that is, whenever the extraction a marginal tonne of waste displaces through recycling is worth more
than collecting, treating and emitting it costs. Since $p^S$ rises as reserves deplete
(Section~\ref{sec:sp:B:longrun}) while $c^W$ and $p^P$ need not, this is not a pathological case but a
natural late-transition one: waste becomes an asset. Two consequences follow, and they are worth
recording because they are a useful internal check on any numerical solution.

First, by~\eqref{eq:sp:B:costate:M} the sign of $p^M$ follows the sign of $p^W$, so
by~\eqref{eq:sp:B:q} the wedge on capital \emph{flips together with} the shadow price of waste: when
waste becomes an asset, embodying material in capital becomes a cost rather than a saving, and $q>1$.
The two flips are simultaneous and share a single cause, so a solution in which they do not coincide is
wrong. Second, any numerical implementation must therefore leave $p^W$, and every instrument built from
it, unbounded in sign.

\smalltitle{The multiplier $\eta$.} By~\eqref{eq:sp:B:foc:cs}, $\eta>0$ only where $\Phi^IG=R$: every
tonne entering production is embodied in new capital, and $\Omega=0$. This is a corner of the
accounting, not of the technology, and it is generic only in an economy investing at close to its entire
material throughput. Where it binds, $\eta$ enters $B^{\mathrm B}$ and hence raises the value of both virgin and
secondary material by the same amount, as it must --- the constraint is on mass, and mass is mass
whatever its source. It also raises $q$ by $\Phi^I\eta$ in~\eqref{eq:sp:B:foc:I}: at the corner, capital
formation is competing for a genuinely scarce tonne. The corner has not been characterised beyond this;
see~\ref{op:sp:corner}.

\smalltitle{The remaining signs.} Collecting~\eqref{eq:sp:B:costate:S}--\eqref{eq:sp:B:costate:M}:
$p^S>0$ (reserves are an asset), $p^X<0$ (cumulative discovery is a liability), $p^P>0$ (pollution is a
liability, so the costate $-\lambda p^P$ is negative as posited), $p^M$ has the sign of $p^W$, and $q$
has the opposite sign of the wedge. Only $p^W$ and $p^M$ are of ambiguous sign, and they are of the same
one.


\subsection{The recycling block}\label{sec:sp:B:recycling}

The two conditions~\eqref{eq:sp:B:foc:varpi} and~\eqref{eq:sp:B:foc:KR} determine how circular the
economy is, and they respond to a single sufficient statistic, $B^{\mathrm B}$.

\smalltitle{Recycling capital.} \eqref{eq:sp:B:foc:KR} reads $F_K=a'(x)B^{\mathrm B}$: capital is allocated so that
its marginal product in production equals its marginal product in recycling, the latter being $a'(x)$
tonnes of secondary material each worth $B^{\mathrm B}$. This is a standard equalisation, but the right-hand side is
not standard: $B^{\mathrm B}$ is a shadow price that depends on the scarcity of the resource, so the recycling
sector's demand for capital inherits the entire dynamics of the reserve.

\smalltitle{The treatment margin.} \eqref{eq:sp:B:foc:varpi} reads
$\alpha B^{\mathrm B}+\left(1-d^W\right)p^P=dc^T/d\varpi$: treatment is pushed to where its marginal cost equals the
value of the material it yields plus the emission it prevents. The two terms are the two distinct
reasons to treat waste, and they are separately identified --- an economy with $d^W=1$ (treatment that
does not immobilise anything) still treats waste, purely to recycle it; an economy with $\alpha=0$ still
treats waste, purely to avoid emissions.

\smalltitle{Circularity rises with scarcity.} The comparative statics of the block are unambiguous,
which is not obvious given that $\alpha$ and $a'$ move in opposite directions.

\begin{proposition}[Scarcity raises circularity]\label{prop:sp:B:circularity}
Hold $F_K$ and $p^P$ fixed and let $B^{\mathrm B}$ rise. Then $x$ rises, $a(x)$ rises, the marginal yield $\alpha$
rises, $\varpi$ rises, and the recirculated share $a\varpi$ rises.
\end{proposition}

\begin{proof}
From~\eqref{eq:sp:B:foc:KR}, $a'(x)=F_K/B^{\mathrm B}$ falls, and $a''<0$ gives $dx/dB>0$; $a'>0$ then gives
$da/dB>0$. For the marginal yield, $\alpha=a(x)-a'(x)x$ has $d\alpha/dx=a'-a''x-a'=-a''x>0$, so $\alpha$
rises with $x$ and hence with $B^{\mathrm B}$. Therefore $\alpha B^{\mathrm B}$ rises, so the left-hand side
of~\eqref{eq:sp:B:foc:varpi} rises; $d^2c^T/d\varpi^2>0$ by~\eqref{eq:sp:setup:waste-cost} then gives
$d\varpi/dB>0$. Both factors of $a\varpi$ rise.
\end{proof}

The mechanism of the paper is this proposition combined with~\eqref{eq:sp:B:arbitrage}
and~\eqref{eq:sp:B:costate:S}. Depletion raises the in-situ rent $p^S$; \eqref{eq:sp:B:arbitrage} passes
the rise one-for-one into $B^{\mathrm B}$; and Proposition~\ref{prop:sp:B:circularity} converts it into a higher
recirculated share. By~\eqref{eq:sp:setup:ledger-solved} the circularity multiplier
$\left(1-a\varpi\right)^{-1}$ rises with it, so a given material throughput $R$ is supported by ever less
virgin extraction. The transition to a circular economy is, in this model, a scarcity-driven
reallocation of capital, not a pollution-driven one --- the pollution channel enters $B^{\mathrm B}$ only through
$d^Wp^P$, which is bounded by the damage function, whereas $p^S$ is not bounded at all.

Note also what does \emph{not} appear in the transition mechanism: no $\omega^j$, no $\Omega$, no
$\Phi^Y$, and, by~\eqref{eq:sp:B:arbitrage}, no $p^W$ either. The date at which the economy turns
circular is a function of extraction costs, resource rents, damages and the recycling technology only.
This is a strong prediction and a useful one, since it means the transition is not sensitive to the
least well identified part of the accounting. It is also a prediction that Version A does not share.


\subsection{Stationarity and the long run}\label{sec:sp:B:longrun}

It is natural to look for a steady state. The model as specified does not have one, and the reason is
instructive.

\begin{proposition}[No interior stationary point]\label{prop:sp:B:nosteadystate}
At any point where all five states of Problem~\ref{prob:sp:B} are stationary, $D=N=0$ and
$W\left(1-a\varpi\right)=0$. Since $a(x)<1$ for all finite $x$, it follows that $W=R=R^R=0$: the economy
uses no materials.
\end{proposition}

\begin{proof}
$\dot X=D=0$ forces $D=0$, and then $\dot S=D-N=0$ forces $N=0$, and hence $\mathcal N=0$. $\dot K=0$
gives $I=0$ and $G=\delta K$; $\dot M^K=0$ then gives $\Phi^IG=\delta M^K$, so the two capital terms
cancel in the ledger~\eqref{eq:sp:B:ledger}, leaving $W=R$. But $R=N+a\varpi W=a\varpi W$, so
$W\left(1-a\varpi\right)=0$. With $\varpi\le1$ and $a<1$ the factor is strictly positive.
\end{proof}

The culprit can be identified exactly, and it is not the closure: the proof uses only the state
equations and the ledger, both of which are common to the two versions. A stationary reserve with
positive extraction requires $D=N>0$, since $\dot S=D-N$; extraction must be replenished by discovery.
But $\dot X=D$ makes cumulative discovery a monotone state, so $D>0$ is itself incompatible with
stationarity. The only route to an interior rest point is therefore closed off by the single assumption
that exploration costs depend on cumulative past discovery. Drop it --- set $C^D_X=0$, so that $X$
ceases to be a state --- and a genuine steady state with $D=N>0$ exists; this is
switch~\ref{sw:sp:CDX} in Section~\ref{sec:sp:schemes}, and it is the cheapest available repair. Note
that removing exploration altogether does \emph{not} repair it: with $\dot S=-N$, stationarity forces
$N=0$ directly. Section~\ref{sec:sp:setup:primitives} already flagged the corresponding structural
feature --- there is no upper bound on $S$, so the resource is never exhausted, only ever more costly ---
and Proposition~\ref{prop:sp:B:nosteadystate} is its dynamic counterpart: with $C^D_X>0$ the economy has
no resting point, it merely becomes asymptotically circular.

That is in fact the substantive content. Read the proposition in the limit rather than at a point. As
$t\to\infty$ the rent $p^S$ grows, $B^{\mathrm B}$ grows with it by~\eqref{eq:sp:B:arbitrage}, and
Proposition~\ref{prop:sp:B:circularity} sends $x\to\infty$, $a\to1$ and $\varpi\to1$. Along such a path
$N\to0$ while $R$ need not: from~\eqref{eq:sp:setup:ledger-solved} with $\dot M^K\to0$ and
$\mathcal N\to0$,
\begin{align}
R &= \frac{N-a\varpi\dot M^K+a\varpi\mathcal N}{1-a\varpi}\;\longrightarrow\;\frac{0}{0},
\label{eq:sp:B:limit}
\end{align}
an indeterminate form whose value is whatever material throughput the recycling sector can sustain. The
emission flow~\eqref{eq:sp:B:Xi} tends to
$\left(1-\varpi+d^W\varpi\right)W-d^Wa\varpi W\to0$, so the limiting economy is materially closed and
emission-free, and it operates at positive scale. This is the sense in which a fully circular economy is
the model's long run: not a steady state that is reached, but a boundary that is approached, and
approached only because extraction becomes infinitely expensive.

\smalltitle{The quasi-stationary block.} Even though the system has no rest point, the price block is
worth recording at $\dot{}=0$, both because it is what a slowly-moving path looks like locally and
because it is the natural initialisation for the numerical solution. Setting all price derivatives to
zero in~\eqref{eq:sp:B:costate:K}--\eqref{eq:sp:B:costate:M} and using $r=\rho$ from $\dot\lambda=0$:
\begin{align}
F_K &= q\left(\rho+\delta\right),
&
p^S &= -\frac{C^N_S}{\rho},
&
p^X &= -\frac{C^D_X}{\rho},
\notag\\
p^P &= \frac{d^{\mathrm B}}{\rho+\theta\!\left(P\right)+\theta'\!\left(P\right)P},
&
p^M &= \frac{\delta p^W}{\rho+\delta},
&
q &= 1-\Phi^Ip^W\frac{\rho}{\rho+\delta},
\label{eq:sp:B:quasistationary}
\end{align}
with $p^W$ from~\eqref{eq:sp:B:foc:W-solved}, and the materially stationary ledger
\begin{align}
N+\mathcal N &= W\left(1-a\varpi\right)
\qquad\text{whenever }\dot M^K=0,
\label{eq:sp:B:material-stationary}
\end{align}
which is~\eqref{eq:sp:setup:ledger-solved} at $\dot M^K=0$, and which requires $\Phi^K=\Phi^I$
by~\eqref{eq:sp:setup:phiK-motion}. Virgin extraction plus displaced mass equals waste not recirculated;
$N\to0$ as $a\varpi\to1$.

\smalltitle{A caveat on the multiplier.} \eqref{eq:sp:B:material-stationary} makes the quantitative
stakes of the circularity multiplier visible. As $a\varpi\to1$ a given $N$ supports an unbounded $W$, so
the model's mass flows become arbitrarily sensitive to the recirculation assumption in exactly the
region the paper is about. The assumption in question is that recycled material re-enters production
contemporaneously --- $R^R=a\varpi W$ at the same instant --- which in continuous time is innocuous but
in a discrete-time implementation ties the number of circuits per period to the period length. It is
flagged here because $R/N$ and $W/N$ should be reported alongside any headline result, and their
sensitivity to the period length checked; see~\ref{op:sp:period}.


\subsection{The complete system}\label{sec:sp:B:system}

For reference, and as the specification the quantitative implementation must reproduce, the optimality
system of Version B consists of:
\begin{itemize}[leftmargin=1.6em,itemsep=0.15em]
\item five state equations~\eqref{eq:sp:B:states-motion} in $K,S,X,P,M^K$;
\item six price equations: \eqref{eq:sp:B:costate:K} for $q$, \eqref{eq:sp:B:costate:S} for $p^S$,
\eqref{eq:sp:B:costate:X} for $p^X$, \eqref{eq:sp:B:costate:P} for $p^P$, \eqref{eq:sp:B:costate:M} for
$p^M$, and \eqref{eq:sp:B:euler} for $C$ (equivalently for $\lambda$);
\item five static conditions determining $I,N,D,\varpi,K^R$: \eqref{eq:sp:B:foc:I},
\eqref{eq:sp:B:foc:N}, \eqref{eq:sp:B:foc:D}, \eqref{eq:sp:B:foc:varpi}, \eqref{eq:sp:B:foc:KR};
\item two definitions of prices: \eqref{eq:sp:B:foc:W} for $p^W$ and~\eqref{eq:sp:B:damage} for $d^{\mathrm B}$;
\item the goods constraint~\eqref{eq:sp:B:goods}, the ledger~\eqref{eq:sp:B:ledger}, and the
complementary-slackness block~\eqref{eq:sp:B:foc:cs};
\item the definitions~\eqref{eq:sp:B:definitions} for $R^R,R,Y,G$ and~\eqref{eq:sp:B:Xi} for $\Xi$;
\item initial conditions on the five states and the transversality conditions~\eqref{eq:sp:B:tvc}.
\end{itemize}
The accounting objects $\Omega$, $\Phi^Y$ and $\Phi^K$ are computed \emph{after} the fact,
from~\eqref{eq:sp:B:Omega} and~\eqref{eq:sp:setup:phiK-motion}; by
Proposition~\ref{prop:sp:B:irrelevance} they do not feed back.

\begin{remark}[The reduced formulation]\label{rem:sp:B:reduced}
Substituting the solved ledger~\eqref{eq:sp:setup:ledger-solved} into the goods constraint and dropping
$W$ from the control set gives a problem in six controls and the same five states, whose first-order
conditions are~\eqref{eq:sp:B:foc:C}--\eqref{eq:sp:B:foc:KR} with $p^W$ replaced everywhere by the
right-hand side of~\eqref{eq:sp:B:foc:W}. The two formulations are equivalent because
\eqref{eq:sp:B:foc:W} is precisely the condition that makes the substitution valid; what is lost is the
price itself. Since $p^W$ is the object the market economy implements through a waste tax, and since the
sign of $p^W$ turns out to determine the sign of two separate instruments simultaneously
(Section~\ref{sec:sp:B:signs}), it is worth keeping.
\end{remark}

\begin{remark}[What is not verified]\label{rem:sp:B:soc}
The conditions above are necessary. Sufficiency requires concavity of the maximised Hamiltonian in the
states, which is not established here: $F$ is concave and $u,-v$ are, but $\mathcal R$ enters through a
constant-returns function of two controls, and the ledger is bilinear in $\varpi$ and $W$. Uniqueness of
the optimum is likewise not established, and it is invoked later when the market implementation is
claimed to reproduce \emph{the} planner's allocation. Both are open; see~\ref{op:sp:soc}.
\end{remark}


\subsection{Comparison with Version A}\label{sec:sp:B:compare}

The two versions differ in one assumption --- whether $\Phi^I$ is exogenous or inherits the economy-wide
intensity $\Omega$ --- and the difference propagates in a way that is worth setting out in one place.
Everything in this subsection refers to Section~\ref{sec:sp:A}, and it is the only part of
Section~\ref{sec:sp:B} that does.

\smalltitle{The whole difference is one scalar.} Define, as in~\eqref{eq:sp:A:zeta},
\begin{align}
\chi &\equiv \frac{\Phi^IG}{R}\in\left[0,1\right),
&
\zeta &= \chi\left(p^W-p^M\right),
&
\Delta &= \zeta\Omega ,
\label{eq:sp:B:compare:Delta}
\end{align}
the share of material input embodied in new capital, and the wedge it generates. Every Version-A
condition is its Version-B counterpart with a unit of final good \emph{spent} on activity $j$ repriced
at $\left(1+\Delta\omega^j\right)$, a unit \emph{produced} repriced at $\left(1-\Delta\omega^Y\right)$,
and the multiplier $\eta$ on the mass constraint replaced by $\zeta$. Setting $\Delta=\zeta=0$ and
reinstating $\eta$ recovers Section~\ref{sec:sp:B:foc} exactly.

\begin{remark}[When the two versions coincide]\label{rem:sp:coincide}
$\Delta=0$ iff $\chi\left(p^W-p^M\right)=0$. The two closures therefore deliver identical allocations in
exactly three circumstances: $\phi^I=0$ (capital embodies no material, so $\chi=0$); $G=0$ (no gross
formation, so nothing is embodied); or $p^W=p^M$ (postponing waste is worth nothing, which
by~\eqref{eq:sp:B:costate:M} requires $r=0$ or $p^W=0$). Otherwise they differ at first order, and the
difference is $O\!\left(\Omega\chi p^W\right)$ --- small when the economy embodies little of its
throughput in durables, large when it embodies much. It is worth reporting $\chi$ alongside $a\varpi$ for
this reason: it is the statistic that says how much the choice of closure can matter.
\end{remark}

\smalltitle{What is common to both.} Four results survive the change of closure unchanged in form, and
they are the results the paper turns on. The materials arbitrage holds in both --- $B^{\mathrm B}=p^{N,\mathrm B}+d^Wp^P$
in~\eqref{eq:sp:B:arbitrage}, $B=p^N+d^Wp^P$ in~\eqref{eq:sp:A:arbitrage} --- and
in both the shadow price of waste cancels from it, so recycling defers waste rather than avoiding it
under either closure. The circularity proposition holds in both, with the same proof, because it uses
only the curvature of $a$ and the convexity of $c^T$. The absence of an interior steady state holds in
both, with the same proof, because it uses only the state equations and the ledger. And capital is a
temporary materials sink in both, with $q$ below its goods cost whenever $p^W>0$. No setting of the
parameters recovers a specification in which recycling avoids waste rather than deferring it; mass
conservation is not a modelling option in this framework.

\smalltitle{Where they differ.} Table~\ref{tab:sp:versions} collects the differences.

\begin{table}[H]
\centering
\caption{Version B and Version A compared.}
\label{tab:sp:versions}
\setlength{\tabcolsep}{5pt}
\renewcommand{\arraystretch}{1.05}
\begin{tabular}{@{}p{0.26\textwidth}p{0.33\textwidth}p{0.33\textwidth}@{}}
\toprule
 & Version B ($\Phi^I$ exogenous) & Version A ($\Phi^I=\Omega\phi^I$) \\
\midrule
Materials restriction & $\Phi^IG\le R$ binds as an inequality, multiplier $\eta\ge0$
& none; $\Omega\ge0$ holds by construction \\
Solving for $W$ & one division, \eqref{eq:sp:setup:ledger-solved}
& positive root of a quadratic, \eqref{eq:sp:setup:ledger-quadratic} \\
$\Omega$ & computed after the fact & solved jointly, a control \\
Relative intensities $\omega^j$ & do not affect the allocation & affect every condition \\
Wedge on spending & none & $1+\Delta\omega^j$ \\
Wedge on output & none & $1-\Delta\omega^Y$ \\
$\lambda$ & equals $u'(C)$ & equals $u'(C)/\left(1+\Delta\omega^C\right)$ \\
Euler equation & standard, \eqref{eq:sp:B:euler} & extra term in $\frac{d}{dt}\ln\left(1+\Delta\omega^C\right)$ \\
Capital price $q$ & $1-\Phi^I\left(p^W-p^M\right)$ & $1-\Phi^I\left(1-\chi\right)\left(p^W-p^M\right)$ \\
$p^W$ & explicit, \eqref{eq:sp:B:foc:W-solved} & implicit, \eqref{eq:sp:A:foc:W-solved} \\
Signs that flip together & $p^W,p^M,q$ & $p^W,p^M,q,\Delta$ \\
\bottomrule
\end{tabular}
\end{table}

\smalltitle{The genuine wedge and the spurious one.} The comparison worth dwelling on is between
Version A's consumption wedge and the artefact of Section~\ref{sec:sp:B:nowedge}. The two have the same
shape. \eqref{eq:sp:B:spurious-C} gives $u'(C)=\lambda\left(1+\Omega^Cp^W\right)$ from a posited
reduced form with fixed intensities; \eqref{eq:sp:A:foc:C} gives
$u'(C)=\lambda\left(1+\Delta\omega^C\right)$ from Version A. Both look like a consumption tax levied on
environmental grounds, and only one of them is real.

They are told apart by the multiplier. The spurious wedge is proportional to $p^W$; the genuine one to
$\chi\left(p^W-p^M\right)$. The difference is exactly the difference in what consuming a unit does. Under
the reduced form it reduces economy-wide waste without reducing the mass drawn in, which conservation of
mass does not permit --- the wedge prices a margin that does not exist. Under Version A total mass is
still conserved, and what consuming a unit does is shift a given mass out of the durable stock and into
the current waste flow; the wedge prices the timing of disposal, not its quantity, which is why it
collapses to zero when timing is worthless ($p^W=p^M$) or when nothing is stored ($\chi=0$). Version B
has no wedge at all because it has no such margin either: there, $\Phi^I$ is fixed, so the split between
stock and flow is not something the composition of spending can move.

\smalltitle{What the choice of closure decides.} Two things, and neither is a robustness check.

The first is what the calibration exercise \emph{is}. Under Version B the relative intensities are a
measurement bridge: they translate the model's mass flows into published material-flow accounts and do
nothing else, and by Proposition~\ref{prop:sp:B:irrelevance} no observation of the allocation identifies
them. Under Version A they are technology, identified in principle from the allocation itself. The
calibration strategy therefore has to be settled jointly with the choice of closure; see~\ref{op:sp:identify}.

The second is how much of Section~\ref{sec:sp:setup:materialaccounting} is load-bearing. Under Version B
that section can be read as bookkeeping, and the process-level detail earns its keep entirely through
what it \emph{removes} --- the spurious wedges. Under Version A the same detail is part of the technology
and determines the allocation directly. Neither closure is adopted as the baseline here, and
Section~\ref{sec:sp:schemes} classifies the available simplifications under both.
